\documentclass[11pt]{article}
\usepackage[margin=1in]{geometry}
\usepackage[T1]{fontenc}
\usepackage{amsmath,booktabs,graphicx,xcolor,enumitem}
\usepackage[hidelinks]{hyperref}
\setlist[itemize]{leftmargin=*,nosep,topsep=3pt}
\setlength{\emergencystretch}{2em}

\newcommand{\todo}[1]{\par\smallskip\noindent
  {\color{blue!60!black}\textit{[TODO: #1]}}\par\smallskip}
\newcommand{\degc}{\ensuremath{{}^{\circ}\mathrm{C}}}

\title{Short-Term Temperature Forecasting with Jena Climate Data}
\author{Kareem Washington \and John Aramis Scott}
\date{CSCI 6840: Data Mining \quad Fall 2026}

\begin{document}
\maketitle

\begin{abstract}
\noindent\textit{Skeleton: replace this outline with one paragraph of
100 to 150 words.}
\begin{itemize}
  \item Problem: does historical weather information improve temperature
  forecasts beyond temperature history alone?
  \item Approach: cleaned hourly Jena observations. 24-hour histories.
  Persistence, linear regression, and random forest at 1, 6, and 24 hours.
  \item Evaluation: chronological training through 2014, validation in 2015,
  and held-out 2016 MAE/RMSE on matched examples.
  \item Finding to summarize: weather inputs reduced RMSE by 5.3 to 9.4\%
  within the same method. Linear regression led at 1/6 hours and the forest
  narrowly led at 24 hours. End with the one-station/one-test-year limitation.
\end{itemize}
\end{abstract}

\section{Introduction and Problem Statement}
\label{sec:introduction}
\textbf{Research question.} Do historical weather measurements improve
1-, 6-, and 24-hour temperature forecasts beyond past temperature alone?
\todo{Write a short motivation paragraph about temperature forecasting for
energy or agricultural planning. These are potential applications.
The project did not measure operational benefits.}
\begin{itemize}
  \item Define the task as supervised time-series regression: forecast a
  future hourly mean temperature in \degc{} using information available at
  the forecast origin.
  \item State the objectives: compare learned methods with persistence.
  Isolate the benefit of weather inputs within each method. Examine how
  errors and large misses change with forecast horizon.
  \item Close with a brief approach/contribution summary: cleaned data,
  lagged features, chronological evaluation, and 15 matched comparisons.
\end{itemize}

\section{Background and Related Work}
\label{sec:background}
\todo{Write one concise paragraph on
Gaddam and Jafari's Jena study~\cite{gaddam}: ARIMA/SARIMA and Box-Jenkins
models, seasonal structure, and external predictors. Verify the final
publication year and manuscript version before submission.}
\begin{itemize}
  \item Position our work as a controlled comparison of lagged linear
  regression and random forest, rather than a reproduction of that study.
  Do not compare its published scores numerically with ours: evaluation
  periods, horizons, and predictor availability differ.
  \item Briefly explain persistence as a transparent benchmark and the
  difference between a weighted linear predictor and an ensemble of trees.
  Connect historical windows and preprocessing to course
  forecasting material~\cite{course}.
\end{itemize}

\section{Methodology}
\label{sec:methods}
\subsection{Data and Preprocessing}
\todo{Turn these verified implementation facts into a compact paragraph.
Cite the Jena data source~\cite{jena}.}
\begin{itemize}
  \item 420,551 raw 10-minute records and 14 weather variables. Remove
  327 exact duplicates before hourly aggregation.
  \item Mask sensor values $\leq -999$. Retain hours with six distinct
  timestamps. Compute right-labeled hourly means and circular mean wind
  direction. Result: 70,034 hourly records, with gaps preserved.
  \item The hourly series ends January 1, 2017. That endpoint is outside
  the 2016 test targets. Fit median imputation on training data only.
  Also fit training-only scaling for linear regression.
\end{itemize}

\subsection{Forecast Targets, Features, and Models}
At origin $t$, the input window contains the latest completed hour and the
previous 23 hours. Each horizon has a separate direct prediction model:
\[
  \widehat{T}_{t+h}=f_h(\mathbf{x}_{t-23},\ldots,\mathbf{x}_{t}),
  \qquad h\in\{1,6,24\}\text{ hours}.
\]
\begin{itemize}
  \item Temperature-only inputs: 24 features. Weather inputs: 336 features
  ($24\times14$), including temperature, pressure, humidity, dew point,
  and wind. Align histories and targets by timestamp across missing hours.
  \item Persistence: $\widehat{T}_{t+h}=T_t$. Learned methods:
  ordinary linear regression and random forest via scikit-learn~\cite{sklearn}.
  \item Forest: 100 trees, minimum leaf size 2, feature fraction 0.33,
  seed 6840. Select depth 12 or 18 by validation RMSE for each input/horizon
  pair. Selected depths were 18 at 1/6 hours and 12 at 24 hours.
\end{itemize}
\todo{Explain why these models and the common 24-hour window make a useful
comparison. Mention Python, pandas, NumPy, Matplotlib, and joblib, and
complete their citations. Focus on the key implementation decisions.}

\section{Experiments and Results}
\label{sec:results}
\subsection{Experimental Setup and Metrics}
\begin{itemize}
  \item Partition chronologically by target time: 2009 to 2014 training,
  2015 validation, 2016 test. Freeze settings before test scoring.
  Final models remain fitted on training data without a validation refit.
  \item At each origin, use observed history then available. This evaluates
  repeated direct forecasts, rather than a year-long recursive forecast.
  \item Five configurations per horizon: persistence plus two learned
  methods with two input sets. All five share the same test timestamps
  within a horizon: $n=8{,}685$, $8{,}675$, and $8{,}639$, respectively.
\end{itemize}
For errors $e_i=\widehat{T}_i-T_i$, report
$\mathrm{MAE}=n^{-1}\sum_i|e_i|$ and
$\mathrm{RMSE}=\sqrt{n^{-1}\sum_i e_i^2}$ in \degc{}. Lower is better.
RMSE is the primary comparison. Optionally report $R^2$.

\subsection{Main Comparisons}
Table~\ref{tab:results} summarizes the test scores.
\todo{Explain what the scores demonstrate: weather helps within both
methods at every horizon. Compare gains over persistence separately from
gains over temperature-only inputs. Calculate percentages before rounding.}

\begin{table}[htbp]
\centering
\small
\setlength{\tabcolsep}{5pt}
\caption{Held-out 2016 errors in \degc{}. Each cell is MAE / RMSE.
Bold marks the lowest observed RMSE at each horizon.
Values are rounded to three decimal places.}
\label{tab:results}
\begin{tabular}{llccc}
\toprule
Method & Inputs & 1 hour & 6 hours & 24 hours \\
\midrule
Persistence & Latest temperature & 0.659 / 0.918 & 3.106 / 4.105 & 2.456 / 3.224 \\
Linear regression & Temperature & 0.400 / 0.581 & 1.699 / 2.230 & 2.459 / 3.157 \\
Linear regression & Weather & \textbf{0.383 / 0.550} & \textbf{1.552 / 2.020} & 2.286 / 2.942 \\
Random forest & Temperature & 0.438 / 0.624 & 1.675 / 2.201 & 2.423 / 3.111 \\
Random forest & Weather & 0.400 / 0.573 & 1.568 / 2.040 & \textbf{2.260 / 2.918} \\
\bottomrule
\end{tabular}
\end{table}

\begin{itemize}
  \item Best observed RMSE reductions from persistence: 40.1\%, 50.8\%,
  and 9.5\% at 1, 6, and 24 hours.
  \item Within-method RMSE reductions from weather inputs: linear
  regression 5.3/9.4/6.8\%. Forest 8.2/7.3/6.2\% in horizon order.
  \item The forest's 24-hour advantage over weather-input linear regression
  is about 0.024\degc{} RMSE. No statistical significance test was performed.
  Temperature-only linear regression slightly worsens 24-hour MAE relative
  to persistence, so avoid claiming every learned model wins on every metric.
\end{itemize}

\subsection{Forecast Example and Error Analysis}
\begin{figure}[htbp]
\centering
\fbox{\parbox[c][0.65in][c]{0.88\linewidth}{\centering
Figure placeholder: observed vs. predicted temperature\\
Insert a 24-hour forecast plot for the first test week.}}

\caption{Planned example: 24-hour forecasts and observations during the
first test week of 2016, plotted at target time. Each prediction uses
information available 24 hours earlier.}
\label{fig:forecast}
\end{figure}
\todo{Discuss Figure~\ref{fig:forecast} after inserting the plot. Explain
the January 7 warming miss (9.16\degc{} for the weather-input forest) and
contrast its 2.26\degc{} MAE with its 5.88\degc{} 95th-percentile absolute
error. The percentile describes observed test errors, not a prediction
interval. The example week illustrates behavior. Annual scores establish
overall accuracy.}

\section{Discussion and Conclusion}
\label{sec:discussion}
\todo{Write one or two concise paragraphs from the following prompts.}
\begin{itemize}
  \item What worked: extra weather history improved both learned methods.
  Linear regression remained competitive. Answer the research question
  within the scope of this experiment.
  \item Challenges/lessons: duplicate records, missing timestamps, sensor
  sentinels, wind angles, and leakage control required explicit decisions.
  Explain one concrete challenge and its implemented solution.
  \item What did not work as well: large day-ahead misses remained.
  Persistence's 24-hour error was lower than its 6-hour error. A daily
  cycle is a possible explanation that was not tested directly.
  \item Limits: one station, one held-out year, one history length, small
  forest search, correlated predictors, no per-variable ablation or
  uncertainty analysis. A small forest advantage does not establish
  significance, and model ranking alone does not prove overfitting.
  \item Future work: seasonal-naive baselines, calendar/circular wind
  features, alternative history lengths, rolling time splits, additional
  years/stations, and uncertainty estimates. Distinguish these proposed
  extensions from completed experiments.
\end{itemize}

\clearpage
\begin{thebibliography}{9}
\bibitem{gaddam}
D. C. Gaddam and R. Jafari,
``Temperature Forecasting on the Jena Climate Dataset Using SARIMA and
Box-Jenkins Models,'' in \textit{Proceedings of the Future Technologies
Conference (FTC) 2025, Volume 3}, pp. 320 to 340.
\href{https://doi.org/10.1007/978-3-032-07995-4_22}{doi:10.1007/978-3-032-07995-4\_22}.

\bibitem{jena}
Max Planck Institute for Biogeochemistry,
\textit{Jena Climate 2009 to 2016}, weather-station observations,
distributed through the TensorFlow/Keras dataset archive.
\url{https://storage.googleapis.com/tensorflow/tf-keras-datasets/jena_climate_2009_2016.csv.zip}.

\bibitem{sklearn}
F. Pedregosa et al., ``Scikit-learn: Machine Learning in Python,''
\textit{Journal of Machine Learning Research}, vol. 12,
pp. 2825 to 2830, 2011.
\url{https://jmlr.org/papers/v12/pedregosa11a.html}.

\bibitem{course}
H. Hedayati, \textit{CSCI 6840 Data Mining: Forecasting and Data
Preprocessing}, course lecture slides, Fall 2026.

\bibitem{software}
\textit{Software citation placeholder.}
\todo{Add official citations for Python, pandas, NumPy, Matplotlib, and
joblib in the same numbered style. Include the software versions used and replace this placeholder entry.}
\end{thebibliography}
\end{document}
